\section{Appendix: Implementation Details}
\label{sec:appendix}

\subsection{Lever Ontology Used in the Pilot}

The bounded candidate set in the current implementation is instantiated
from the following ontology:
\begin{quote}
\small
graffiti removal; litter removal; facade repair; storefront
transparency increase; localized greenery addition; lighting repair;
signage decluttering; crosswalk repainting; lane marking repainting;
surface cleaning; shutter repair; tree canopy management.
\end{quote}

\subsection{Prompt Contracts}

The code-level planner, editor, and critic prompts are implemented in
the repository configuration and enforce a structured contract rather
than free-form prompting.

\noindent\textbf{Planner contract.}
The planner is instructed to return a bounded set of single-lever
candidates grounded in visible scene elements. Each candidate must
contain exactly:
\begin{quote}
\small
\texttt{lever\_concept}, \texttt{scene\_support},
\texttt{target\_object}, \texttt{intervention\_direction},
\texttt{edit\_template}, and \texttt{edit\_plan}.
\end{quote}
Additional hard constraints require ontology membership, one lever per
candidate, locality, plausibility, and prompt-only-friendly edits.

\noindent\textbf{Edit prompt contract.}
The image-edit prompt explicitly forbids global restyling, relighting,
camera changes, background changes, and adding or removing non-target
objects. The prompt requires the model to preserve the input scene and
modify only the specified target object according to the lever plan.

\noindent\textbf{Critic contract.}
The automated critic receives the original image, the edited image, and
the intended lever specification, and returns structured JSON for four
binary checks:
\begin{quote}
\small
\texttt{same\_place\_preserved}, \texttt{is\_localized},
\texttt{is\_realistic}, and \texttt{is\_plausible},
\end{quote}
plus brief repair guidance. An edit is accepted only if all four checks
are true.

\subsection{Current Pilot Configuration}

The reproducible pilot reported in the paper uses:
\begin{itemize}
    \item SPECS subset size $N{=}20$;
    \item five cities with 4 images per city;
    \item per-city stratification by baseline safety score and visual
    complexity;
    \item candidate budget $K{=}3$;
    \item stochastic budget $T{=}1$ per fixed lever;
    \item \texttt{flux-kontext-max} for prompt-only editing;
    \item ViT-PP2 for auxiliary safety scoring.
\end{itemize}

The corresponding driver command is:
\begin{quote}
\small
\texttt{uv run python scripts/run\_repro\_pipeline.py --prefix specs\_repro\_n20 --n-total 20 --candidate-budget 3 --max-attempts 1 --score-device mps}
\end{quote}

\subsection{Exported Artifacts}

The main paper-facing run writes:
\begin{itemize}
    \item a candidate-level CSV with full lever identities and critic
    fields;
    \item a scored CSV with baseline score, edited score, and auxiliary
    delta;
    \item a per-image auxiliary summary CSV;
    \item generated paper figures derived directly from the scored
    outputs.
\end{itemize}

Human-evaluation export and ingestion paths are implemented in the
repository but are optional in the current draft and were not used to
produce the reported pilot results.
